\begin{center}
    \section*{\fontsize{20}{20}\selectfont Chapter 6}
\end{center}
\vspace{10mm}

\section{Conclusion}

This thesis shows that graph neural networks provide a viable and efficient way to segment brain tumors, challenging the assumption that you need volumetric 3D convolutions for competitive medical image analysis. The framework we proposed achieves 92.92\% Dice coefficient through 5-fold ensemble methods on BraTS 2021 while delivering 6.9$\times$ faster inference and 156$\times$ fewer parameters compared to 3D U-Net. This enables deployment on consumer-grade hardware and could help democratize access to AI-assisted diagnostics.

Key contributions? First, superpixel-based graph construction that exploits tumor sparsity for 890$\times$ dimensionality reduction. Second, validation that a 5-layer GraphSAGE architecture is optimal through ablation studies. Third, discovering batch size sensitivity in medical imaging GNN tasks. And fourth, rigorous validation with patient-level stratified cross-validation and 15 integrity checks addressing data leakage concerns.

Going back to the research question—\textit{Can we achieve competitive accuracy with practical efficiency gains for clinical deployment?}—this thesis provides an affirmative answer. As healthcare AI scales globally, efficiency-aware architectures become essential for equitable healthcare access.


\subsection{Limitations}

Several limitations are worth acknowledging. First, there's preprocessing overhead of 2-3 seconds per patient that impacts real-time workflows. Second, handcrafted features require domain expertise and might not generalize to all tumor types. Third, we focused on binary segmentation rather than multi-class delineation. Fourth, validation is limited to BraTS 2021 gliomas—a single dataset. Fifth, there's still a 1-2\% performance gap below transformer methods. And sixth, clinical integration barriers including regulatory approval, hospital IT integration, and real-world validation with radiologists remain unaddressed.


\subsection{Future Works and Direction}

Several promising directions emerge from this work. Architectural enhancements through attention mechanisms and hybrid CNN-GNN designs could push accuracy higher. Multi-task extensions for multi-class segmentation and survival prediction would increase clinical utility. Preprocessing optimization via differentiable clustering and GPU acceleration could reduce that 2-3 second overhead. 

For clinical translation, we'd need radiologist studies, regulatory approval, and multi-institutional trials. Broader applications to other anatomical regions and longitudinal tumor tracking also seem promising.

This thesis establishes graph neural networks as a compelling paradigm for efficient brain tumor segmentation. It demonstrates that competitive medical image analysis doesn't require massive volumetric models. As healthcare AI scales globally, efficiency-aware architectures will be critical for ensuring advanced diagnostics reach diverse populations regardless of resource availability.